\section{Sharpness at every exponent}
\label{sec:sharpness}

The constant-modulus equality cases of
\cref{thm:mass-energy-sandwich} already show sharpness in finite
dimension.  We now match the arithmetic normalization and show that
the exponent transfers cannot be improved from only
\[
 S_X\ll Q_X,\qquad
 L_X\ll X^{o(1)}N_{0,X}/Q_X.
\]
Rounding an integer support size changes only fixed or subpower
factors and will be suppressed.

\begin{proposition}[Sharpness of mass-to-energy loss]
\label{prop:sharp-mass-energy}
Fix \(\mu\ge0\).  There are finite sequences with
\[
 S_X\asymp Q_X,\qquad
 L_X\le\frac{N_{0,X}}{Q_X},
\]
for which
\[
 M_X=X^{-\mu}N_{0,X},
\qquad
 D_X\asymp
 X^{-2\mu}\frac{N_{0,X}^2}{Q_X}.
\]
Thus the exponent \(2\mu\) in
\eqref{eq:energy-from-mass} is best possible.
\end{proposition}

\begin{proof}
Choose \(S_X=\lceil Q_X\rceil\) nonzero coordinates and give each
coordinate modulus
\[
 t_X=\frac{X^{-\mu}N_{0,X}}{S_X}.
\]
Then
\[
 M_X=S_Xt_X=X^{-\mu}N_{0,X}
\]
and
\[
 D_X=S_Xt_X^2
=\frac{X^{-2\mu}N_{0,X}^2}{S_X}
\asymp
X^{-2\mu}\frac{N_{0,X}^2}{Q_X}.
\]
Also \(t_X\ll N_{0,X}/Q_X\).  Arbitrary unit phases may be attached
to the coordinates.
\end{proof}

\begin{proposition}[Sharpness of energy-to-mass and support losses]
\label{prop:sharp-energy-mass}
Fix \(\lambda\ge0\) in a range for which
\(X^{-\lambda}Q_X\to\infty\).  There are finite sequences satisfying
\[
 L_X=\frac{N_{0,X}}{Q_X}
\]
and
\begin{align}
 S_X&\asymp X^{-\lambda}Q_X,\label{eq:sharp-support}\\
 M_X&\asymp X^{-\lambda}N_{0,X},\label{eq:sharp-mass}\\
 D_X&\asymp
 X^{-\lambda}\frac{N_{0,X}^2}{Q_X}.
\label{eq:sharp-energy}
\end{align}
Consequently, neither the mass exponent \(\lambda\) nor the support
exponent \(\lambda\) in
\cref{thm:two-sided-transfer} can be improved from the stated
envelope.
\end{proposition}

\begin{proof}
Choose
\[
 S_X=\lfloor X^{-\lambda}Q_X\rfloor
\]
coordinates, all of modulus \(N_{0,X}/Q_X\).  Then
\eqref{eq:sharp-support}--\eqref{eq:sharp-energy} follow by direct
calculation.  This sequence has constant modulus on its support and
therefore attains equality on both sides of
\eqref{eq:mass-energy-sandwich}.
\end{proof}

\begin{remark}[Very sparse range]
If \(X^{-\lambda}Q_X\) is bounded, the lower support conclusion is
at most a constant-scale statement.  The theorem remains valid, but
there is no stronger power-scale support conclusion to extract from
the envelope alone.
\end{remark}

\begin{example}[Why the support upper bound is needed]
Let \(M=1\) and spread this mass uniformly over \(K\) coordinates.
Then \(D=1/K\).  Allowing \(K\) to exceed the physical
\(O(Q_X)\) determinant range would make energy arbitrarily smaller
than the scale in \eqref{eq:energy-from-mass}.  Thus the first
transfer uses the verified physical support interval, not merely
finiteness.
\end{example}

\begin{example}[Why the singleton bound is needed]
A one-point sequence of amplitude \(T\) has
\[
 D=T^2,\qquad M=T,\qquad L=T.
\]
Without an upper bound for \(L\) in the physical normalization, an
energy lower bound does not determine the normalized mass scale
\(N_{0,X}\).  The reverse transfer therefore genuinely uses the
TPC-32 singleton estimate.
\end{example}
